%=============================================================================
% AGENT 10: COMPLETE ENGINEERING SYNTHESIS
% Integration of All Agent Outputs into a Unified Specification
% for the DFD Psi-Bubble Propulsion Device
%
% Sources: Agent 1 (freefall derivation, kappa-lambda independence),
%          Agent 2 (missile capture, psi-wave momentum flux),
%          Agent 3 (Noether loophole, Q-amplification),
%          Original 20-agent campaign (psi_bubble_propulsion.tex)
%=============================================================================

\documentclass[12pt,aps,prd,superscriptaddress,nofootinbib,longbibliography]{revtex4-2}
\usepackage{amsmath,amssymb,amsfonts,amsthm}
\usepackage{bm,physics}
\usepackage{siunitx}
\usepackage{booktabs}
\usepackage{tcolorbox}
\usepackage{graphicx}
\usepackage{enumitem}

\newcommand{\avec}{\mathbf{a}}
\newcommand{\neff}{n_{\mathrm{eff}}}
\newcommand{\psib}{\psi_{\mathrm{b}}}
\newcommand{\ka}{k_a}
\newcommand{\etac}{\eta_c}
\newcommand{\Rb}{R_{\mathrm{b}}}
\newcommand{\Rc}{R_{\mathrm{c}}}
\newcommand{\Lshell}{\ell_{\mathrm{shell}}}
\newcommand{\calB}{\mathcal{B}}

\newtheorem{theorem}{Theorem}
\newtheorem{corollary}{Corollary}

\begin{document}

\title{The $\psi$-Bubble Propulsion Device:\\
Complete Engineering Synthesis from DFD First Principles}

\author{DFD Research Programme --- Integration Agent}
\date{28 March 2026}

\begin{abstract}
We synthesize all research outputs from the DFD $\psi$-bubble
investigation into a single, honest engineering specification.  Three
independent propulsion mechanisms are identified, analyzed, and ranked:
(A)~the Q-amplified cross-term force in an external gravitational field,
(B)~the local $n_{\mathrm{eff}}$ modification creating a self-generated
gravity well, and (C)~$\psi$-wave radiation recoil.  Mechanism~C is
shown to be gravitationally suppressed ($G/c^5$ factor) and negligible.
Mechanism~A, operating through the Noether loophole (ship-in-external-field
is an open system), yields kilonewton-class thrust for achievable
parameters ($B = 20$~T, $Q \sim 4 \times 10^9$, $V = 10$~m$^3$).
Mechanism~B provides the dominant propulsion channel once above
threshold: the $\kappa(\eta - \eta_c)$ modification to $n_{\mathrm{eff}}$
creates an asymmetric refractive-index distribution, and all matter
(ship, occupants, atmosphere) follows geodesics of the optical metric in
freefall.  We present: complete device specifications, the above-threshold
bubble profile, the $\psi$-gradient and resulting acceleration, the
self-force mechanism, all UAP-matching observable signatures, time dilation
effects, the momentum budget (with honest caveats), power requirements,
and the critical experimental test.
\end{abstract}

\maketitle
\tableofcontents

%=============================================================================
\section{The Three Propulsion Mechanisms}\label{sec:three-mechanisms}
%=============================================================================

The DFD formalism, through the agents' analyses, reveals three distinct
mechanisms by which a ferrite-superconductor shell could generate thrust.
Their physical origins and scaling laws are fundamentally different.

\subsection{Mechanism A: Q-Amplified Cross-Term in External Gravity}

\paragraph{Physics (Agent 3).}
The DFD master equation contains a nonlinear $\ka a^2/c^2$ term.
Expanding the acceleration field into external and internal components,
$\avec = \avec_{\mathrm{ext}} + \avec_{\mathrm{ship}}$, the cross-term
$2\avec_{\mathrm{ext}} \cdot \avec_{\mathrm{ship}}$ couples the ship's
EM energy to the external gravitational gradient.  The resulting force is:
\begin{equation}
F_{\mathrm{cross}} = \ka \frac{U_{\mathrm{EM}}}{c^2} \, g \, \Gamma,
\qquad \ka = \frac{3}{8\alpha} \approx 51.4,
\label{eq:F-cross}
\end{equation}
where $g$ is the local gravitational acceleration and $\Gamma$ is a
geometric overlap factor.

\paragraph{Resonant amplification.}
When the EM field is modulated at the $\psi$-mode resonant frequency
$\Omega_\psi = \pi c/R$ (Agent 3, Theorem 1), the $\psi$-mode amplitude
is enhanced by the quality factor $Q$:
\begin{equation}
\boxed{F_Q = Q \cdot \ka \cdot \frac{U_{\mathrm{EM}}}{c^2} \cdot g \cdot \Gamma.}
\label{eq:F-Q}
\end{equation}

\paragraph{Noether legitimacy (Agent 3, \S1).}
This does not violate Noether's theorem: the ship is immersed in an
external $\psi$-field ($\psi_{\mathrm{ext}} = 2gz/c^2$) that explicitly
breaks translational symmetry.  The device is a gravitational sail,
pushing off the external field.  In free space ($g \to 0$), this force
vanishes identically.

\paragraph{Numerical assessment.}
For $B = 20$~T, $V = 10$~m$^3$, $Q = 4 \times 10^9$, $\Gamma = 0.3$,
near Earth ($g = 9.8$~m/s$^2$):
\begin{equation}
F_Q = 4 \times 10^9 \times 51.4 \times \frac{1.6 \times 10^9}{(3 \times 10^8)^2}
\times 9.8 \times 0.3 \approx 9.8~\mathrm{kN}.
\end{equation}
This equals the weight of a 1000~kg capsule---sufficient for $1g$ hover.


\subsection{Mechanism B: Local $n_{\mathrm{eff}}$ Modification (Self-Generated Gravity Well)}

\paragraph{Physics (Agent 1, \S2; Agent 2, \S1).}
Above the threshold $\eta > \etac = \alpha/4$, the effective refractive
index is:
\begin{equation}
\neff = \exp\!\Big[\psi + \kappa(\eta - \etac)\,\Theta(\eta - \etac)\Big],
\label{eq:neff}
\end{equation}
where $\eta = U_{\mathrm{EM}}/(\rho c^2)$.  An asymmetric EM distribution
(stronger on one hemisphere) creates $\nabla\neff \neq 0$.  All matter
follows geodesics of the optical metric:
\begin{equation}
\avec = \frac{c^2}{2}\nabla\psi_{\mathrm{eff}}
= \frac{c^2}{2}\frac{\nabla\neff}{\neff}.
\label{eq:geodesic}
\end{equation}

\paragraph{Key distinction from Mechanism A.}
This is NOT gravitational radiation or cross-term coupling.  The
$\neff$ modification is a \emph{local, adiabatic} change to the optical
metric---a constitutive modification of the medium itself.  Agent 1's
kappa-lambda independence analysis proves that this channel operates through
the $\kappa = 51.4$ enhancement, independent of the $\lambda$ channel
that governs $\psi$-wave generation (which is suppressed by
$|\lambda - 1| \lesssim 3 \times 10^{-5}$).

\paragraph{Acceleration at threshold.}
At threshold ($\eta = \etac$) with asymmetric drive of fractional
asymmetry $\Delta\eta/\etac$, the acceleration is:
\begin{equation}
a = \frac{c^2}{2}\,\kappa\,\frac{\Delta\eta}{R_{\mathrm{shell}}}.
\end{equation}
For $\kappa = 1$, $\Delta\eta = \etac = 1.82 \times 10^{-3}$,
$R = 5$~m:
\begin{equation}
a = \frac{(3 \times 10^8)^2}{2} \times 1 \times
\frac{1.82 \times 10^{-3}}{5} = 1.6 \times 10^{13}~\mathrm{m/s^2}.
\end{equation}

This is an \emph{enormous} acceleration---far more than needed.  The
actual operating point would be controlled by modulating $\Delta\eta$
to the desired level, with the threshold as an on/off switch.

\paragraph{Critical caveat.}
Agent 1 identifies a serious concern: the $\psi$-gradient varies across
the shell, and the ship must self-consistently ``fall into'' its own
generated well.  The acceleration depends on the gradient at the
\emph{ship's} location, not at the shell boundary.  For a well-designed
bubble where the gradient is smooth on the shell scale, the ship
accelerates toward the higher-$\neff$ hemisphere.  But the detailed
self-consistent solution requires specifying how the asymmetric $\neff$
distribution maps to the gradient at the center of mass.


\subsection{Mechanism C: $\psi$-Wave Radiation Recoil}

\paragraph{Physics (Agent 2, psi-wave momentum flux).}
A pulsed above-threshold EM source radiates $\psi$-waves through the
gravitational wave equation $\Box\psi = -(8\pi G/c^4)\mathcal{T}$.
The geometric asymmetry of the source creates a net momentum flux.

\paragraph{Fatal suppression.}
The radiated power scales as (Agent 2, Eq.~5.9):
\begin{equation}
P_\psi \approx 2.6 \times 10^{-26}\,\kappa^2~\mathrm{W},
\end{equation}
and the thrust as:
\begin{equation}
F_{\mathrm{Mode~II}} \approx 3 \times 10^{-35}~\mathrm{N}.
\end{equation}
The $G/c^5 = 2.75 \times 10^{-53}$~s/J factor kills this channel.
All $\psi$-wave radiation from laboratory sources is gravitationally
suppressed to unmeasurable levels.

\begin{tcolorbox}[colback=red!5!white, colframe=red!50!black,
  title=Mechanism C Verdict]
$\psi$-wave radiation recoil is \textbf{not viable} for propulsion.
The $G/c^5$ suppression is a fundamental barrier.  This mechanism is
ruled out.
\end{tcolorbox}


\subsection{Resolution: Which Mechanism Dominates?}

\begin{tcolorbox}[colback=green!5!white, colframe=green!50!black,
  title=The Answer to the Critical Question]
\textbf{Mechanisms A and B are complementary.  Mechanism B dominates
for UAP-like performance.  Mechanism A provides the near-Earth
self-force.}

\smallskip
\textbf{Near a gravitating body:} Both A and B operate.  Mechanism A
(cross-term force) provides $F_Q = Q \ka (U_{\mathrm{EM}}/c^2) g \Gamma$,
which scales with $g$.  Mechanism B ($\neff$ modification) creates a
local gravity well that is independent of external $g$.

\smallskip
\textbf{In free space:} Mechanism A shuts off ($g \to 0$).
Mechanism B still operates---the device creates its own gravity well and
falls into it.  However, the momentum budget becomes the critical
question (see Section~\ref{sec:momentum}).

\smallskip
\textbf{For UAP-like performance:} The observed signatures (100$g$+
acceleration, no sonic boom, trans-medium, object capture) require
Mechanism B.  The $\psi$-gradient must extend far enough into the
surrounding medium to entrain matter (air, water) in freefall.
\end{tcolorbox}


%=============================================================================
\section{Complete Device Specification}\label{sec:device-spec}
%=============================================================================

\subsection{Architecture}

\begin{enumerate}
\item \textbf{Outer shell: Ferrite ($\epsilon \approx \mu$).}
NiZn ferrite tiles (50~mm $\times$ 50~mm $\times$ 10~mm), impedance-matched
to vacuum.  This maximizes EM energy storage density for given input
power.  NiZn selected for operation above 100~MHz (high resistivity,
loss tangent $\sim 10^{-2}$ at 1~GHz).

\item \textbf{Inner shell: Type-II superconductor.}
Nb$_3$Sn on Cu substrate (50~$\mu$m coating), $T_c = 18.3$~K,
$H_{c2} = 24$~T.  Meissner shielding blocks all EM fields from interior.
Creates sharp boundary: $\eta > \etac$ in ferrite, $\eta = 0$ inside.

\item \textbf{Structural shell: Ti-6Al-4V.}
5~mm thick, provides mechanical support and thermal conduction.
Yield strength 880~MPa at cryogenic temperature.  Safety factor $> 14$
for EM radiation pressure at $U_0 = 10$~MJ.
\end{enumerate}

\subsection{Dimensions}

\begin{table}[h]
\caption{Baseline device dimensions.}
\begin{ruledtabular}
\begin{tabular}{lcc}
Parameter & Baseline & Flight \\
\hline
Geometry & Oblate spheroid & Prolate spheroid \\
Semi-major axis & 3.0~m & 5.0~m \\
Semi-minor axis & 1.5~m & 3.0~m \\
Ferrite thickness & 10~mm & 10~mm \\
SC thickness & 2~mm (Nb) / 50~$\mu$m (Nb$_3$Sn) & 50~$\mu$m \\
Structural shell & 5~mm Ti & 5~mm Ti \\
Shell mass & 4,920~kg & $\sim$10,000~kg \\
\end{tabular}
\end{ruledtabular}
\end{table}

\subsection{Operating Parameters}

\begin{table}[h]
\caption{Operating parameters for the Q-amplified cross-term (Mechanism A).}
\begin{ruledtabular}
\begin{tabular}{llll}
Parameter & Symbol & Value & Status \\
\hline
Self-coupling & $\ka$ & 51.4 & DFD prediction (unverified) \\
Threshold & $\etac$ & $\alpha/4 \approx 1.82 \times 10^{-3}$ & DFD prediction \\
Shell radius & $R$ & 3--5~m & Engineering choice \\
Magnetic field & $B$ & 10--20~T & Achievable (SC magnets) \\
Shell volume & $V$ & 10--100~m$^3$ & Engineering choice \\
Quality factor & $Q$ & $4 \times 10^9$ & SC cavity range \\
$\psi$-mode frequency & $f_\psi$ & 50~MHz ($R = 3$~m) & VHF band \\
Geometric factor & $\Gamma$ & 0.3 (conservative) & Design-dependent \\
EM rotation freq. & $f_{\mathrm{rot}}$ & 25~MHz & Half of $f_\psi$ \\
\end{tabular}
\end{ruledtabular}
\end{table}

\begin{table}[h]
\caption{Threshold conditions (Mechanism B).}
\begin{ruledtabular}
\begin{tabular}{lll}
Medium & $\rho$~[kg/m$^3$] & $B_c$ to reach $\etac$ \\
\hline
Ultra-high vacuum ($10^{-8}$) & $10^{-8}$ & 0.64~T \\
Near-vacuum ($10^{-4}$) & $10^{-4}$ & 64~T \\
Sea-level air & 1.2 & $7 \times 10^3$~T \\
Ferrite shell & 5000 & $4.5 \times 10^5$~T \\
\end{tabular}
\end{ruledtabular}
\end{table}

\paragraph{Critical observation.}
The threshold is easily crossed in vacuum ($B_c = 0.64$~T for
$\rho = 10^{-8}$~kg/m$^3$), moderately achievable in near-vacuum,
but impossible in dense matter at achievable field strengths.  The
above-threshold region necessarily exists \emph{outside} the shell,
in the surrounding vacuum or low-density gas, not within the dense
ferrite shell itself.

This is the design-critical insight: the ferrite shell stores and shapes
the EM field, but the propulsively active $n_{\mathrm{eff}}$-modified
region is the vacuum/low-density gas \emph{surrounding} the shell where
$\eta > \etac$.


%=============================================================================
\section{The Above-Threshold Region Profile}\label{sec:bubble-profile}
%=============================================================================

\subsection{Extent of the Bubble}

The above-threshold region is defined by $\eta(\mathbf{x}) > \etac$,
i.e., the region where the local EM energy density exceeds
$\etac \rho c^2$.  For a magnetic field decaying from the shell:

\paragraph{Dipole decay.}
$B(r) \approx B_0 (R/r)^3$ for $r > R$.  The above-threshold boundary
is at:
\begin{equation}
r_{\mathrm{AT}} = R \left(\frac{B_0}{B_c}\right)^{1/3}.
\end{equation}

For $B_0 = 20$~T, $R = 5$~m, in near-vacuum ($B_c = 0.64$~T):
\begin{equation}
r_{\mathrm{AT}} = 5 \times (20/0.64)^{1/3} = 5 \times 3.15 = 15.7~\mathrm{m}.
\end{equation}

The bubble extends $\sim 3R$ from the shell surface.

\paragraph{For the sech$^2$ model (Agent 2).}
With bubble radius $\Rb$ and shell thickness $\Lshell$, the profile is:
\begin{equation}
\psi_{\mathrm{eng}}(r) = \psib \, \operatorname{sech}^2\!\left(
\frac{r - \Rb}{\Lshell}\right),
\end{equation}
with maximum gradient:
\begin{equation}
|\nabla\psi|_{\mathrm{max}} \approx \frac{\psib}{\Lshell}.
\end{equation}

\subsection{Capture Radius (Agent 2)}

The capture radius---where the $\psi$-gradient produces acceleration
exceeding a threshold $a_{\mathrm{thresh}}$---is:
\begin{equation}
\Rc = \Rb + \frac{\Lshell}{2}\ln\!\left(
\frac{4\,a_{\mathrm{max}}}{a_{\mathrm{thresh}}}\right).
\end{equation}

For $a_{\mathrm{max}} = 100g$, $\Rb = 50$~m, $\Lshell = 10$~m:
\begin{itemize}
\item At $a_{\mathrm{thresh}} = 1g$: $\Rc \approx 80$~m.
\item At $a_{\mathrm{thresh}} = 0.01g$: $\Rc \approx 103$~m.
\end{itemize}

The capture zone extends roughly 30--50~m beyond the nominal bubble
shell.  This is a logarithmic extension---doubling $a_{\mathrm{max}}$
adds only $\sim 3.5$~m.


%=============================================================================
\section{The $\psi_{\mathrm{eff}}$ Gradient and Acceleration}
\label{sec:gradient}
%=============================================================================

\subsection{Required Gradient for $1g$}

From $\avec = (c^2/2)\nabla\psi$:
\begin{equation}
|\nabla\psi|_{1g} = \frac{2g}{c^2} = 2.18 \times 10^{-16}~\mathrm{m}^{-1}.
\end{equation}

This is identical to Earth's $\psi$-gradient at the surface---as it must
be, since $1g$ acceleration requires the same $\psi$-gradient regardless
of source.

\subsection{Maximum Acceleration from the Bubble}

For the sech$^2$ profile:
\begin{equation}
a_{\mathrm{max}} = \frac{c^2}{2}\frac{\psib}{\Lshell}.
\end{equation}

\paragraph{$\psib$ for $100g$ ($\sim 10^3$~m/s$^2$) with $\Lshell = 10$~m:}
\begin{equation}
\psib = \frac{2 a_{\mathrm{max}} \Lshell}{c^2}
= \frac{2 \times 10^3 \times 10}{(3 \times 10^8)^2}
\approx 2.2 \times 10^{-13}.
\end{equation}

This is a minuscule $\psi$-perturbation producing macroscopic acceleration.

\subsection{Acceleration Profile}

The acceleration as a function of distance from bubble center:
\begin{equation}
a(r) = \frac{c^2\psib}{\Lshell}\,\operatorname{sech}^2\!\left(
\frac{r - \Rb}{\Lshell}\right)
\left|\tanh\!\left(\frac{r - \Rb}{\Lshell}\right)\right|.
\end{equation}

\begin{itemize}
\item At $r = \Rb$ (shell peak): $a = 0$ (gradient vanishes at the
maximum of a sech$^2$).
\item At $r = \Rb \pm \Lshell/2$: $a \approx a_{\mathrm{max}}$.
\item Far field ($r \gg \Rb$): $a \sim 4 a_{\mathrm{max}}
\exp[-2(r - \Rb)/\Lshell]$ (exponential decay).
\end{itemize}


%=============================================================================
\section{The Self-Force Mechanism}\label{sec:self-force}
%=============================================================================

\subsection{The Central Question}

How does the device accelerate itself?  The previous analysis showed
the five-channel force budget.  We now resolve which channels provide
useful thrust.

\subsection{Channel Assessment}

\begin{table}[h]
\caption{Force channels ranked by magnitude.}
\begin{ruledtabular}
\begin{tabular}{llll}
Channel & Mechanism & Magnitude & Viable? \\
\hline
A: Cross-term & $Q\ka(U/c^2)g\Gamma$ & $\sim$~kN (near Earth) & \textbf{Yes} \\
B: $\neff$ modification & $(c^2/2)\kappa\Delta\eta/R$ & $\sim 10^{13}$~m/s$^2$
  & \textbf{Yes}$^*$ \\
C: $\psi$-wave recoil & $G/c^5$ suppressed & $\sim 10^{-35}$~N & No \\
D: EM radiation & Photon momentum $P/c$ & $\sim 10^{-3}$~N & Negligible \\
E: Plasma sling & Gravitational scattering & Environment-dependent & Supplementary \\
\end{tabular}
\end{ruledtabular}
\begin{flushleft}
$^*$See critical discussion below regarding self-consistency.
\end{flushleft}
\end{table}

\subsection{How Mechanism A Works: The Gravitational Sail}

Near a massive body, the cross-term force (Eq.~\ref{eq:F-Q}) acts as a
gravitational sail.  The Q-enhanced $\psi$-mode inside the shell creates
a rectified DC $\psi$-gradient through the nonlinear $a^2$ term.  This
gradient, interacting with the external $\nabla\psi_{\mathrm{ext}}$,
produces a net force.

\textbf{The recoil is carried by the external gravitational field.}
Just as a solar sail pushes off sunlight, the $\psi$-bubble pushes off
Earth's (or any body's) gravitational $\psi$-field.  Momentum flows
through the $\psi$-field to the source mass.

\subsection{How Mechanism B Works: Falling Into Your Own Well}

The device creates an asymmetric $\neff$ distribution outside the shell.
In the DFD framework, higher $\neff$ corresponds to a deeper gravitational
potential.  The ship ``falls toward'' the higher-$\neff$ region---this is
geodesic motion on the optical metric, identical in principle to
gravitational freefall.

\paragraph{The self-consistency question.}
Can an object fall into its own gravitational well?  In Newtonian gravity,
a uniform shell creates no field inside itself (shell theorem).  But the
$\psi$-bubble is:
\begin{enumerate}
\item \textbf{Asymmetric}: one hemisphere has higher $\neff$ than the other.
\item \textbf{External}: the above-threshold region is outside the shell,
not inside.
\item \textbf{Continuously regenerated}: the EM field is continuously
driven, maintaining the asymmetry as the ship moves.
\end{enumerate}

The shell theorem does not apply because (a) the source is not spherically
symmetric, and (b) the $\neff$ modification is in the surrounding medium,
not within the shell.  The ship at the center of the shell experiences a
net gradient from the asymmetric external $\neff$ distribution.

\paragraph{The ambient medium channel.}
When operating in atmosphere or plasma, the above-threshold region extends
into the surrounding gas.  Air molecules in the gradient region accelerate
at $(c^2/2)\nabla\psi$ toward the higher-$\neff$ side.  This is
gravitational attraction of the ambient medium, and the reaction force
on the ship provides thrust.  The air effectively serves as the reaction
mass---not expelled mechanically, but gravitationally accelerated.

This is the Channel~5 / Mode~I mechanism operating at the scale of the
ambient medium.  It is equivalent to creating a gravity well and letting
the surrounding matter fall in asymmetrically.

\subsection{Combined Operation}

Near Earth:
\begin{itemize}
\item Mechanism A (cross-term): $F_Q \sim 10$~kN, proportional to $g$.
\item Mechanism B ($\neff$ well): Dominant once above threshold;
controlled by $\Delta\eta$ asymmetry.
\item Channel E (ambient medium): Air/plasma drawn into the well provides
reaction force.
\end{itemize}

In deep space:
\begin{itemize}
\item Mechanism A: Fades as $g \to 0$.
\item Mechanism B: Still operates---the local $\neff$ modification is
independent of external gravity.
\item Momentum budget: Carried by the $\psi$-field deformation extending
to infinity (see Section~\ref{sec:momentum}).
\end{itemize}


%=============================================================================
\section{Observable Signatures Matching UAP Reports}\label{sec:uap}
%=============================================================================

\subsection{No Sonic Boom (Agent 2, \S6)}

\paragraph{Physical mechanism.}
A conventional aircraft pushes air aside mechanically, creating compression
waves.  The $\psi$-bubble does not push air---it creates a $\psi$-gradient
that accelerates air molecules gravitationally.

\paragraph{Two regimes.}
\begin{itemize}
\item \textbf{Entraining} ($v < v_{\mathrm{capture}} = c\sqrt{\psib/2}$):
Air is fully captured and co-moves with the bubble.  The bubble carries its
own atmosphere.  No relative motion between bubble surface and air, hence
no shock.

\item \textbf{Transparent} ($v \gg v_{\mathrm{capture}}$): Air passes
through the gradient with negligible deflection (energy transfer per
molecule $\propto \psib^2/v^2$, which is tiny for $\psib \sim 10^{-13}$).
The bubble is invisible to the atmosphere.  No shock.
\end{itemize}

In both regimes: \textbf{no sonic boom, no thermal wake, no aerodynamic
drag signature.}

For $\psib \sim 10^{-13}$ (100$g$ bubble), $v_{\mathrm{capture}} \approx 200$~m/s.
Below 200~m/s, air is entrained; above 200~m/s, air passes through.  At
no point is a shock wave generated.

\subsection{Instantaneous Direction Changes}

\paragraph{Physical mechanism.}
The ship is in freefall inside the bubble.  Changing the asymmetry
direction of the EM field (electronic thrust vectoring) changes
$\nabla\psi_{\mathrm{eff}}$ instantaneously.  The ship immediately begins
falling in the new direction.

\paragraph{No g-forces.}
By the Weak Equivalence Principle (Agent 1, Theorem 2), all matter inside
the Meissner-shielded interior accelerates identically.  The occupants
experience zero proper acceleration regardless of the coordinate
acceleration.  A $100g$ direction change is experienced as perfect
weightlessness throughout.

\paragraph{Rate limit.}
The electronic reconfiguration time is limited only by the RF phase
control bandwidth ($\sim 100$~kHz), allowing direction changes within
$\sim 10~\mu$s---appearing instantaneous to any external observer.

\subsection{Trans-Medium Travel (Air to Water)}

\paragraph{Physical mechanism.}
The $\psi$-gradient acts identically on all matter regardless of
composition or phase.  Air molecules, water molecules, and the ship
all follow the same geodesic.

Entering water: water molecules in the leading gradient zone accelerate
gravitationally, opening a path.  The bubble transitions seamlessly
because the propulsion mechanism does not depend on the medium's
mechanical properties---only on the $\psi$-gradient.

\paragraph{No splash, no cavitation.}
Water is drawn into the gradient gravitationally, not displaced
mechanically.  The transition produces no splash, no shockwave,
and no cavitation noise.

\subsection{EM Interference on Nearby Electronics}

\paragraph{Physical mechanism.}
The rotating EM field at $\sim 50$~MHz (VHF band) is partially
confined by the shell but inevitably leaks.  The near-field of a
$10^7$--$10^9$~J stored-energy system produces strong EM fields at
VHF/UHF frequencies within several hundred meters.

\paragraph{Expected effects.}
\begin{itemize}
\item Radio interference: strong.  The drive frequency (25--50~MHz) and
harmonics produce broadband RF noise.
\item Vehicle electronics disruption: disrupted within the near-field
zone ($r < \lambda/2\pi \sim 1$~m for the fundamental; harmonics
extend further).
\item Compass deflection: the DC magnetic field from the superconducting
coils (20~T at the shell) decays as $r^{-3}$ but is detectable at
substantial distances.
\end{itemize}

\subsection{Luminous Glow}

\paragraph{Physical mechanism.}
Multiple contributing effects:
\begin{enumerate}
\item \textbf{Corona discharge}: Strong near-field EM ionizes
surrounding air, producing visible glow (similar to St.\ Elmo's fire).
\item \textbf{Cherenkov-like radiation}: In the $\neff > 1$ region, the
phase velocity of light is $c/\neff < c$.  Charged particles in the
ambient medium moving at $v > c/\neff$ produce Cherenkov radiation.
For $\neff = 1 + \delta n$ with $\delta n \sim 10^{-13}$, this requires
$v > c(1 - 10^{-13})$---essentially cosmic ray speeds, so this effect
is negligible at low $\psib$.
\item \textbf{Plasma recombination}: Ionized air (from RF heating)
recombining produces characteristic emission spectra.
\end{enumerate}

\subsection{Object Capture (Agent 2, Missile Capture Analysis)}

\paragraph{Physical mechanism.}
Any object entering the $\psi$-gradient zone accelerates at
$(c^2/2)\nabla\psi$---the same acceleration as the ship (WEP).
The object is captured in freefall.

\paragraph{Capture dynamics.}
\begin{enumerate}
\item Object enters gradient zone at relative velocity $v_{\mathrm{rel}}$.
\item The gradient accelerates the object in the same direction as the
ship.
\item Inside the bubble ($r < \Rb$), the field is approximately uniform;
the object floats with zero relative acceleration.
\item From outside, the object appears to be ``grabbed'' and carried along.
\end{enumerate}

\paragraph{Capture velocity.}
Full capture requires $v_{\mathrm{rel}} < v_{\mathrm{capture}}
= c\sqrt{\psib/2}$.  For $\psib \sim 2 \times 10^{-11}$
($a_{\mathrm{max}} \sim 3 \times 10^4\,g$): $v_{\mathrm{capture}}
\sim 1000$~m/s (Mach 3 missile capture).

\paragraph{Tidal forces at boundary.}
$\Delta a_{\mathrm{tidal}} = (2 a_{\mathrm{max}}/\Lshell) \times L$,
where $L$ is the object size.  For a 5~m missile at $a_{\mathrm{max}}
= 100g$, $\Lshell = 10$~m: $\Delta a_{\mathrm{tidal}} = 100g$.
Significant but survivable for solid objects.  A thicker shell
($\Lshell = 100$~m) reduces tidal forces 100-fold.

\subsection{No Visible Exhaust}

\paragraph{Physical mechanism.}
The device does not expel mass.  It creates a $\psi$-gradient and falls
into it.  The reaction momentum is carried by:
\begin{itemize}
\item Deformation of the external $\psi$-field (invisible).
\item Gravitational acceleration of ambient medium (air flowing smoothly,
not explosively expelled).
\end{itemize}

There is no exhaust plume, no thermal signature from propellant
combustion, and no reaction mass stream.


%=============================================================================
\section{Time Dilation Effects}\label{sec:time-dilation}
%=============================================================================

\subsection{General Formula (Agent 1, Theorem 3)}

From the physical metric:
\begin{equation}
\frac{d\tau}{dt} = e^{-\psi_{\mathrm{eff}}/2}.
\end{equation}

The fractional time dilation is:
\begin{equation}
\frac{\delta t}{t} = 1 - e^{-\Delta\psi_{\mathrm{eff}}/2}
\approx \frac{\Delta\psi_{\mathrm{eff}}}{2}
\quad (\text{for } |\Delta\psi_{\mathrm{eff}}| \ll 1).
\end{equation}

\subsection{Numerical Estimates}

\begin{table}[h]
\caption{Time dilation for different operating regimes.}
\begin{ruledtabular}
\begin{tabular}{llll}
Regime & $\Delta\psi_{\mathrm{eff}}$ & $\delta t/t$ & Comparison \\
\hline
$1g$ acceleration, $R = 5$~m & $1.1 \times 10^{-15}$ & $5 \times 10^{-16}$
  & Atomic clock precision \\
$100g$, $R = 5$~m & $1.1 \times 10^{-13}$ & $5 \times 10^{-14}$
  & 100$\times$ GPS satellites \\
Strong above-threshold & $1.6 \times 10^{-2}$ & $8 \times 10^{-3}$
  & $10^7 \times$ Earth gravity \\
  ($\eta \sim 10\,\etac$) & & & \\
\end{tabular}
\end{ruledtabular}
\end{table}

The time dilation effect is detectable by atomic clocks for moderate
accelerations and becomes macroscopic ($\sim 1\%$) at strong above-threshold
operation.


%=============================================================================
\section{The Momentum Budget}\label{sec:momentum}
%=============================================================================

\begin{tcolorbox}[colback=yellow!5!white, colframe=yellow!50!black,
  title=Honesty Box: What Is Resolved and What Is Not]

\textbf{Resolved:}
\begin{enumerate}
\item The system is open: the ship couples to the external $\psi$-field
and ambient medium.  Noether's theorem permits $dp_{\mathrm{ship}}/dt \neq 0$.
\item Near a massive body, the cross-term (Mechanism A) transfers momentum
to the Earth through the $\psi$-field.  The recoil carrier is identified:
$\Delta p_\oplus = -\Delta p_{\mathrm{ship}}$.
\item In atmosphere/plasma, the ambient medium provides reaction mass
through gravitational deflection (Channel E).
\item Total momentum is conserved:
$\mathbf{p}_{\mathrm{ship}} + \mathbf{p}_{\psi} + \mathbf{p}_{\mathrm{medium}}
+ \mathbf{p}_{\mathrm{source~mass}} = \mathrm{const}$.
\end{enumerate}

\textbf{Not fully resolved:}
\begin{enumerate}
\item \textbf{Self-consistent Mechanism B in vacuum.}  The local $\neff$
modification creates a gradient that accelerates the ship.  But can the
ship's own EM field pull itself forward?  This is analogous to the
Baron Munchausen bootstrap problem.  The resolution via the $\psi$-field
as momentum carrier requires a quantitative calculation of the $\psi$-field
momentum flux $T^{0i}_\psi = -(c^4/(4\pi G))\dot\psi\nabla\psi$ for
the specific device geometry.  This has not been done to full
self-consistency.

\item \textbf{Free-space operation without external $g$.}
Mechanism A vanishes when $g \to 0$.  Mechanism B may still work if the
$\psi$-field can carry sufficient momentum.  But the $\psi$-field
momentum density scales as $c^4/G$ (large) times $|\nabla\psi|^2$
(tiny), and whether the net flux is nonzero depends on the global
boundary conditions.

\item \textbf{The $Q^2$ scaling claim.}
Agent 3 suggests that in the strongly nonlinear regime, rectification
may give $Q^2$ rather than $Q$ scaling.  This requires a full
nonlinear analysis that has not been performed.

\item \textbf{The $\ka = 51.4$ value.}
The entire force calculation depends on $\ka = 3/(8\alpha)$.  This is a
DFD prediction, not an experimental result.  If $\ka = 0$, both
Mechanisms A and B collapse.
\end{enumerate}
\end{tcolorbox}


%=============================================================================
\section{Power Requirements and Energy Source}\label{sec:power}
%=============================================================================

\subsection{Power Budget}

\begin{table}[h]
\caption{Power budget for the $\psi$-bubble device.}
\begin{ruledtabular}
\begin{tabular}{lcc}
Subsystem & Power (kW) & Notes \\
\hline
RF drive (maintain cavity) & 43 & $P = 2\pi f U_0/Q$, $U_0 = 10$~MJ \\
Cryogenics (2~K) & 100 & 4 GM cryocoolers at 25~kW each \\
Control and avionics & 5 & FPGA, sensors, comms \\
Life support & 10 & 6 crew, recycling \\
Margin (20\%) & 32 & \\
\hline
\textbf{Total} & \textbf{190} & \\
\end{tabular}
\end{ruledtabular}
\end{table}

\subsection{Energy Source}

\begin{itemize}
\item \textbf{Ground testing:} Grid power or diesel generator.  200~kW
sustained.
\item \textbf{Flight system:} Compact fission reactor.  USNC Micro Modular
or equivalent.  5~MW$_{\mathrm{th}}$, 1~MW$_e$ (Brayton cycle, $\eta = 0.2$).
TRISO HALEU fuel, 20-year lifetime.  Total power system mass: 6,000~kg.
\item \textbf{Alternative:} Radioisotope thermoelectric generator (RTG)
for unmanned probes requiring $< 10$~kW.
\end{itemize}

\subsection{Energy Efficiency}

The power to maintain the cavity ($\sim 43$~kW) is independent of the
thrust produced.  This makes the device extraordinarily efficient at
high thrust: a 10~kN thrust at 43~kW input gives an effective
``specific power'' of 233~N/kW, compared to $\sim 0.06$~N/kW for ion
drives.

However, the kinetic power $P_{\mathrm{kinetic}} = Fv$ grows linearly
with velocity.  At $v = 1$~km/s, a 10~kN thrust requires
$P_{\mathrm{kinetic}} = 10$~MW---exceeding the reactor output.  The
energy to accelerate to high velocity comes from somewhere.  In
Mechanism A, it comes from the gravitational potential energy of the
Earth-ship system.  In Mechanism B, the energy source is less clear
and is part of the unresolved momentum budget question.

\subsection{The Critical Energy Question}

\begin{tcolorbox}[colback=yellow!5!white, colframe=yellow!50!black,
  title=Unresolved: Energy Source for Kinetic Energy]
At $1g$ sustained acceleration, after 10 minutes ($v \approx 6$~km/s),
the kinetic energy of a 22,680~kg vehicle is:
\begin{equation}
E_{\mathrm{kin}} = \frac{1}{2}mv^2 = \frac{1}{2} \times 22680 \times 6000^2
\approx 4 \times 10^{11}~\mathrm{J} = 400~\mathrm{GJ}.
\end{equation}
The reactor produces $\sim 10^6$~J/s.  In 10 minutes it delivers
$\sim 6 \times 10^8$~J.  There is a $\sim 700\times$ energy deficit.

For Mechanism A: the energy comes from the gravitational PE of the
Earth-ship system.  As the ship rises, it extracts energy from the
gravitational field at rate $\dot{E} = F_Q \cdot v$.  This is consistent
with energy conservation.

For Mechanism B in free space: the energy source must be identified.
It cannot come from the onboard reactor (insufficient).  It cannot come
from the gravitational field (none nearby).  This is a genuine open
problem.
\end{tcolorbox}


%=============================================================================
\section{The Experimental Test}\label{sec:experiment}
%=============================================================================

\subsection{What Validates the Threshold Physics}

The entire propulsion concept rests on a single testable prediction:
\textbf{above the threshold $\eta > \etac = \alpha/4$, the effective
gravitational coupling of EM energy is enhanced by $\ka = 51.4$.}

\subsection{The Channel 4 Weight Anomaly Experiment}

\begin{table}[h]
\caption{Threshold validation experiment.}
\begin{ruledtabular}
\begin{tabular}{ll}
Parameter & Specification \\
\hline
Solenoid & 10~T, 10~cm bore, SC coil \\
Ambient density & Ultra-high vacuum ($10^{-8}$~kg/m$^3$) \\
$\eta$ at center & $B^2/(2\mu_0 \rho c^2) = 440$ \\
$\eta/\etac$ & $\sim 2.4 \times 10^5$ (far above threshold) \\
Expected $\Delta g$ & $\ka (U_{\mathrm{EM}}/c^2) g / m_{\mathrm{system}}$ \\
& $= 51.4 \times (4 \times 10^4)/(9 \times 10^{16}) \times 9.8$ \\
& $\approx 2.2 \times 10^{-10}$~m/s$^2$ \\
Sensor & SC gravimeter ($10^{-12}$~$g$ sensitivity) \\
Signature & Weight anomaly appearing/vanishing with field \\
& cycling, scaling as $B^2$, with threshold behavior \\
\end{tabular}
\end{ruledtabular}
\end{table}

\paragraph{Go/No-Go.}
Detection of a weight anomaly at the predicted level ($\sim 10^{-11}$~$g$)
confirms the $\ka$-enhancement.  Non-detection to $10^{-12}$~$g$
constrains $\ka < 5$, effectively ruling out the propulsion concept.

\subsection{The $\lambda$ Measurement at SQMS}

An independent test measures the back-reaction coefficient $\lambda$ using
a high-$Q$ SRF cavity in a TE+TM superposition mode (breaking geometry
transparency).  The $\psi$-oscillation at $2\omega$ is detected by
Fabry-Perot interferometry or atom interferometry.

Sensitivity: $|\lambda - 1| > 3 \times 10^{-14}$.

This test validates the $\psi$-wave generation channel ($\lambda$), which
is independent of but complementary to the $\ka$ test.

\subsection{Phase Sequence}

\begin{enumerate}
\item \textbf{Phase 0}: Channel 4 weight anomaly with SC solenoid
in vacuum.  Cost: \$1M.  Timeline: 6 months.
\item \textbf{Phase 1}: $\lambda$ measurement at SQMS.
Cost: \$2M.  Timeline: 12 months.
\item \textbf{Phase 2}: Ferrite-SC coupon test with rotating mode.
Cost: \$5M.  Timeline: 18 months.
\item \textbf{Phase 3}: Subscale demonstrator (0.5~m shell).
Cost: \$20M.  Timeline: 3 years.
\item \textbf{Phase 4}: Full-scale prototype.
Cost: \$200M.  Timeline: 5 years.
\end{enumerate}

The critical path runs through Phase 0.  If the weight anomaly is not
detected, the program stops.


%=============================================================================
\section{Comprehensive Parameter Summary}\label{sec:summary}
%=============================================================================

\begin{table*}[t]
\caption{Complete device specification summary.}
\begin{ruledtabular}
\begin{tabular}{llll}
Category & Parameter & Value & Status \\
\hline
\multicolumn{4}{c}{\textit{Fundamental DFD Parameters}} \\
Self-coupling & $\ka = 3/(8\alpha)$ & 51.4 & Predicted, unverified \\
EM threshold & $\etac = \alpha/4$ & $1.82 \times 10^{-3}$ & Predicted, unverified \\
$\kappa$-$\lambda$ independence & Proven from action & --- & Theoretical (Agent 1) \\
\hline
\multicolumn{4}{c}{\textit{Shell Parameters}} \\
Geometry & Prolate spheroid & $a = 5$~m, $b = 3$~m & Baseline \\
Ferrite & NiZn, $\mu_r \sim 100$ & 10~mm tiles & Existing material \\
Superconductor & Nb$_3$Sn & 50~$\mu$m on Cu & Existing technology \\
Structure & Ti-6Al-4V & 5~mm & Standard aerospace \\
Shell mass & --- & $\sim 5000$--$10000$~kg & Engineering \\
\hline
\multicolumn{4}{c}{\textit{EM Operating Parameters}} \\
Magnetic field & $B$ & 10--20~T & Achievable \\
Stored energy & $U_0$ & 10~MJ & Routine for SC magnets \\
$\psi$-mode frequency & $f_\psi = c/(2R)$ & 50~MHz ($R = 3$~m) & VHF band \\
EM rotation frequency & $f_{\mathrm{rot}}$ & 25~MHz & Half of $f_\psi$ \\
Quality factor (SC cavity) & $Q$ & $10^9$--$10^{11}$ & Demonstrated at SQMS \\
Drive power & $P$ & 43~kW & At $Q = 10^9$ \\
\hline
\multicolumn{4}{c}{\textit{Propulsion Performance (Mechanism A, near Earth)}} \\
Cross-term force ($Q = 4 \times 10^9$) & $F_Q$ & 9.8~kN & For 1000~kg at $1g$ \\
Required $Q$ for $1g$ (1000~kg) & $Q_{1g}$ & $3.7 \times 10^9$ & Achievable \\
Acceleration (22,680~kg vehicle) & $a$ & 0.43~m/s$^2$ & At $F_Q = 9.8$~kN \\
Geometric factor & $\Gamma$ & 0.3 & Conservative \\
\hline
\multicolumn{4}{c}{\textit{Bubble Properties}} \\
Above-threshold radius & $r_{\mathrm{AT}}$ & $\sim 3R$ (15~m) & In near-vacuum \\
Capture radius ($1g$ threshold) & $\Rc$ & $\sim 80$~m & $a_{\mathrm{max}} = 100g$ \\
Capture velocity & $v_c$ & 200~m/s ($\psib \sim 10^{-13}$) & For 100$g$ bubble \\
Tidal force at boundary & $\Delta a$ & $(2a_{\mathrm{max}}/\Lshell) L$ &
  Survivable for thick shell \\
\hline
\multicolumn{4}{c}{\textit{Observable Signatures}} \\
EM transparency & $\delta n$ & $\sim 10^{-13}$ & Undetectable optically \\
Sonic boom & --- & None & Both regimes \\
Time dilation ($1g$) & $\delta t/t$ & $\sim 5 \times 10^{-16}$ & Atomic clock precision \\
EM interference range & --- & $\sim 100$~m & VHF near-field \\
\hline
\multicolumn{4}{c}{\textit{Power System}} \\
Total electrical power & --- & 190~kW & Including margin \\
Energy source & Fission reactor & 1~MW$_e$ & TRISO HALEU, 20~yr \\
Total vehicle mass & --- & 22,680~kg & With crew for 180 days \\
\end{tabular}
\end{ruledtabular}
\end{table*}


%=============================================================================
\section{Honest Assessment and Open Questions}\label{sec:assessment}
%=============================================================================

\subsection{What the Analysis Establishes}

\begin{enumerate}
\item \textbf{The DFD formalism provides a self-consistent framework}
for propulsion via $\psi$-gradient generation.  The math works: geodesic
motion on the optical metric, WEP freefall, threshold activation, and
cross-term forces are all derivable from the DFD postulates.

\item \textbf{The $\kappa$ and $\lambda$ channels are independent}
(Agent 1 proof).  The $\lambda$ suppression ($|\lambda-1| < 3 \times 10^{-5}$)
does not kill the $\kappa = 51.4$ channel.  This is the twelve-order-of-magnitude
window that makes propulsion conceivable.

\item \textbf{The Q-amplified cross-term (Mechanism A) gives kilonewton-class
forces} for achievable parameters near Earth.  The physics is analogous to
a gravitational sail and satisfies energy/momentum conservation by coupling
to the external gravitational field.

\item \textbf{All UAP signatures emerge naturally} from a $\psi$-gradient
bubble: no sonic boom, zero g-forces, trans-medium capability, object
capture by freefall, EM interference from the drive system, and no visible
exhaust.

\item \textbf{$\psi$-wave radiation (Mode II) is not viable.}  The
$G/c^5$ suppression makes laboratory-scale $\psi$-wave thrust negligible
($\sim 10^{-35}$~N).
\end{enumerate}

\subsection{What Remains Unresolved}

\begin{enumerate}
\item \textbf{DFD itself is unverified.}  The entire analysis assumes DFD
is correct---specifically that $\ka = 51.4$ and $\etac = \alpha/4$.
These are predictions of a theory that has not been experimentally tested
at the level required.  The weight anomaly experiment (Phase 0) is the
critical test.

\item \textbf{Mechanism B self-consistency in free space.}
The local $\neff$ modification channel produces enormous accelerations
once above threshold ($\sim 10^{13}$~m/s$^2$).  But the question of
whether a device can self-consistently ``fall into its own well'' without
external gravity has not been answered with a complete self-consistent
calculation.  The shell theorem analogy suggests it might work for
asymmetric geometries, but this requires numerical solution of the full
nonlinear field equation with the device as source.

\item \textbf{Energy budget for sustained acceleration.}
At $1g$ for 10 minutes, the kinetic energy exceeds the reactor output
by $700\times$.  For Mechanism A, this energy comes from gravitational PE.
For free-space operation via Mechanism B, the energy source is unidentified.

\item \textbf{The Q-factor operating in the $\psi$-sector.}
The $Q \sim 10^9$--$10^{11}$ values cited are for the EM cavity mode.
The relevant Q is for the \emph{$\psi$-mode} of the shell, which is
a different quantity that depends on how $\psi$-perturbations dissipate.
This Q has never been measured and may be much lower.

\item \textbf{Materials at operating conditions.}
NiZn ferrite at cryogenic temperatures in 20~T fields: magnetic saturation,
loss tangent changes, and mechanical properties need experimental
characterization.  The $\epsilon \approx \mu$ matching condition must
hold under operating conditions.

\item \textbf{Quantum coherence enhancement ($N_{\mathrm{Cooper}}^p$).}
The coherence exponent $p$ is unknown.  The range $p = 0$ (no enhancement)
to $p = 1$ (full coherence) spans 27 orders of magnitude.  This
uncertainty dominates all performance predictions that depend on the
$\lambda$ channel.
\end{enumerate}

\subsection{The Bottom Line}

The DFD $\psi$-bubble concept is \textbf{internally consistent within
the DFD framework} and \textbf{produces performance predictions matching
UAP observational reports} if the theory's key predictions ($\ka = 51.4$,
$\etac = \alpha/4$) are correct.

The concept is \textbf{experimentally testable}.  A superconducting
solenoid in ultra-high vacuum, monitored by a superconducting gravimeter,
can confirm or rule out the $\ka$-enhancement within 6 months of
approval.  This is the Phase 0 experiment that gates everything else.

If the weight anomaly is detected at the predicted level, the path
to a propulsion demonstrator is clear and uses existing technology
(SC magnets, ferrite materials, RF electronics, compact fission
reactors).  If it is not detected, the concept is falsified.

Either way, the experiment is worth doing.

\end{document}
